\documentclass[11pt]{article}

\usepackage[a4paper, hmargin=2cm, vmargin=2cm]{geometry}
\usepackage[T1]{fontenc}
\usepackage[utf8]{inputenc}
\usepackage[french]{babel}
\usepackage{amsmath, amssymb}
\usepackage{enumitem}
\usepackage{tikz}
\usepackage{pstricks,pst-plot,pst-node}

\pagestyle{empty}

\begin{document}

\begin{center}
\Large \textbf{Devoir n\textsuperscript{o} 3 : Nombre dérivé (55 min)}\\[2mm]
\large 1\textsuperscript{re} spécialité
\end{center}

\medskip

\emph{Tout doit être justifié avec clarté sauf mention explicite contraire. La calculatrice n'est pas autorisée.}

\setlength{\columnseprule}{.5pt}
\setlength{\columnsep}{30pt}

\setcounter{section}{0}
\setcounter{equation}{0}

%==================== Exercice 1 ====================
\noindent\textbf{Exercice 1 }

Sur le graphique ci-dessous, les droites $\Gamma_{f;0}$ et $\Gamma_{f;3}$ sont les tangentes à la courbe $\mathcal{C}_f$ aux points d'abscisses $0$ et $3$.

\begin{center}
\psset{algebraic=true, plotpoints=500, linewidth=1.2pt}
\begin{pspicture}(-3,-4)(6,6)
  % Grille et axes
  \psgrid[subgriddiv=1,gridlabels=7pt,gridcolor=lightgray](-3,-4)(6,6)
  \psaxes[labels=all,ticks=all,dx=1,dy=1,linewidth=1.1pt]{->}(0,0)(-3,-4)(6,6)

  %--- Début du "clipping" (découpe) ---
  \psclip{\psframe(-3,-4)(6,6)}

    % Courbe : f(x) = -(2/27)x^3 + x + 1
    \psplot[linecolor=blue]{-3}{6}{-(2.0/27.0)*x^3 + x + 1}

    % Tangentes
    \psplot[linecolor=red,linewidth=1.2pt]{-3}{6}{x + 1}          % pente +1
    \psplot[linecolor=orange,linewidth=1.2pt]{-3}{6}{-x + 5}       % pente -1

  \endpsclip
  %--- Fin du "clipping" ---

  % Points et étiquettes
  \psdots[dotsize=3pt](0,1)
  \uput[225](0,1){\footnotesize $(0;1)$}
  \psdots[dotsize=3pt](3,2)
  \uput[45](3,2){\footnotesize $(3;2)$}

  \uput[180](-2.2,1){\color{blue}\small $\mathcal{C}_f$}
  \uput[45](2,3.5){\color{red}\footnotesize $\Gamma_{f;0}$}
  \uput[135](4.3,1.5){\color{black}\footnotesize $\Gamma_{f;3}$}

  % Petits segments de pente (repères)
  \psline[linecolor=red,linestyle=dashed](0,1)(1,2)
  \psline[linecolor=orange,linestyle=dashed](3,2)(4,1)

  % Cadre
  \psframe[linewidth=1pt](-3,-4)(6,6)
\end{pspicture}
\end{center}

\begin{enumerate}
\item Compléter par lecture graphique et \textbf{sans justifier} :
  \begin{enumerate}[label=$\bullet$,itemsep=10pt]
    \item $f(0)=\,$\dotfill
    \item $f(3)=\,$\dotfill
    \item $f'(0)=\,$\dotfill
    \item $f'(3)=\,$\dotfill
  \end{enumerate}

\item Déterminer les équations réduites de $\Gamma_{f;0}$ et de $\Gamma_{f;3}$

\end{enumerate}

\bigskip

%==================== Exercice 2 ====================
\noindent\textbf{Exercice 2 }

Soit $f$ la fonction définie sur $\mathbb{R}$ par $f(x)=2x^2-3x$.

\begin{enumerate}
\item Soit $h$ un réel non nul. Calculer et simplifier le taux d'accroissement
\[ \tau_{f;3;3+h}=\frac{f(3+h)-f(3)}{h}. \]

\item Déterminer $\displaystyle \lim_{h\to 0} \tau_{f;3;3+h}$. Que peut-on conclure ?

\item On admet que $f'(2)=5$. Déterminer l'équation de la tangente $\mathcal{D}$ à la courbe $\mathcal{C}_f$ au point $A(2;f(2))$.

\item Le point $K(3;7)$ est-il un point de $\mathcal{D}$ ?

\item Le point $R(-2;-10)$ est-il un point de $\mathcal{D}$ ?
\end{enumerate}

\newpage

%==================== Exercice 3 ====================
\noindent\textbf{Exercice 3 }

Soit $f(x)=\dfrac{1}{x}$, la fonction inverse, définie pour $x > 0$. Soit $a>0$ et $h\neq 0$ tel que $a+h>0$ .

\begin{enumerate}
\item Écrire et simplifier le taux d'accroissement $\tau_{f;a;a+h}$. Montrer que
\[
\tau_{f;a;a+h}=-\frac{1}{a(a+h)}.
\]

\item En déduire que $f$ est dérivable en $a$ et donner le nombre dérivé de $f$ en $a$.

\item En déduire les nombres dérivés $f'(1)$; $f'(2)$;$f'(10)$. 
\end{enumerate}

\end{document}
